\chapter{智慧水利概述与平台架构基础}

\begin{learningobjectives}
通过本章学习，学生应能够：
\begin{enumerate}
    \item 理解智慧水利的基本概念、发展背景及其在水利现代化中的重要作用
    \item 掌握智慧水利平台的总体架构体系，包括感知层、传输层、平台层、应用层的功能和关系
    \item 分析国内外智慧水利发展现状与技术趋势，了解智慧水利平台开发的技术路径
    \item 认识软件工程方法在智慧水利平台开发中的应用价值和实践意义
\end{enumerate}
\end{learningobjectives}

\section{引言}

智慧水利是新一代信息技术与传统水利工程深度融合的产物，代表着水利行业数字化转型的重要发展方向。随着物联网（Internet of Things, IoT）、大数据（Big Data）、人工智能（Artificial Intelligence, AI）、云计算（Cloud Computing）等技术的快速发展，传统的水利管理模式正在发生深刻变革。智慧水利平台作为这一变革的核心载体，集成了现代信息技术与水利专业知识，为水资源管理、水安全保障、水环境保护等提供了全新的解决方案。

软件工程作为一门应用计算机科学、数学和工程原则来设计、开发、维护和测试软件的学科，为智慧水利平台的开发提供了科学的方法论和技术保障。通过系统化、规范化和量化的软件工程方法，可以确保智慧水利平台的高质量、可靠性和可维护性，有效应对复杂水利系统的开发挑战。

\section{第一节
智慧水利平台概述与发展背景}\label{ux7b2cux4e00ux8282-ux667aux6167ux6c34ux5229ux5e73ux53f0ux6982ux8ff0ux4e0eux53d1ux5c55ux80ccux666f}

\subsection{1.1
智慧水利的定义与内涵}\label{ux667aux6167ux6c34ux5229ux7684ux5b9aux4e49ux4e0eux5185ux6db5}

智慧水利是在新一轮科技革命和产业变革背景下，运用物联网、大数据、人工智能、云计算、5G通信、数字孪生等新一代信息技术，与水利工程、水利管理、水利服务深度融合，构建具有预报、预警、预演、预案功能的现代化水利体系。智慧水利通过全面感知、泛在互联、智能分析、协同决策，实现水利业务的数字化转型和智能化升级，提升水利治理体系和治理能力现代化水平。

智慧水利的核心内涵体现在以下几个方面：

\subsubsection{1.1.1 全面感知}\label{ux5168ux9762ux611fux77e5}

通过部署各类传感器、监测设备、卫星遥感、无人机等现代化感知手段，实现对水情、雨情、工情、灾情等水利要素的全方位、全天候、实时动态监测。感知网络覆盖江河湖库、水利工程、用水户等各个环节，形成''天空地''一体化的立体感知体系。

\textbf{主要技术特征：} -
\textbf{多源感知}：集成水文站网、气象站网、工程监测、视频监控、卫星遥感等多种数据源
- \textbf{实时监测}：支持秒级、分钟级数据采集和传输，确保信息的时效性 -
\textbf{智能感知}：运用人工智能技术，实现异常状态的自动识别和预警

\subsubsection{1.1.2 泛在互联}\label{ux6cdbux5728ux4e92ux8054}

构建高速、可靠、安全的水利通信网络，实现感知设备、业务系统、管理部门之间的无缝连接和信息共享。通过5G、光纤、卫星通信等多种通信手段，确保数据传输的实时性和可靠性，打破信息孤岛，促进数据流通和业务协同。

\textbf{网络架构特点：} -
\textbf{多网融合}：整合有线、无线、卫星等多种通信方式 -
\textbf{边缘计算}：在数据源头部署边缘计算设备，提升数据处理效率 -
\textbf{安全传输}：采用加密技术和安全协议，确保数据传输安全

\subsubsection{1.1.3 智能分析}\label{ux667aux80fdux5206ux6790}

运用大数据分析、机器学习、深度学习等人工智能技术，对海量水利数据进行深度挖掘和智能分析，发现数据背后的规律和趋势，为水利决策提供科学依据。通过构建各类水利模型和算法，实现预报预警、趋势分析、风险评估等智能化应用。

\textbf{分析能力包括：} -
\textbf{多维分析}：支持时空多维度的数据分析和可视化展示 -
\textbf{预测建模}：基于历史数据建立预测模型，实现精准预报 -
\textbf{智能识别}：自动识别异常模式和潜在风险

\subsubsection{1.1.4 协同决策}\label{ux534fux540cux51b3ux7b56}

基于智能分析结果，结合专业知识和管理经验，为水利管理和工程调度提供科学决策支持。通过数字仿真、情景推演等技术手段，评估不同决策方案的效果，优化决策过程，提升决策的科学性和有效性。

\textbf{决策支持特点：} -
\textbf{多方案比较}：生成多种调度方案并进行效果评估 -
\textbf{实时优化}：根据实时监测数据动态调整决策方案 -
\textbf{智能推荐}：基于历史经验和模型计算推荐最优决策

\subsection{1.2
智慧水利与传统水利的区别}\label{ux667aux6167ux6c34ux5229ux4e0eux4f20ux7edfux6c34ux5229ux7684ux533aux522b}

传统水利管理主要依靠人工观测、经验判断和定性分析，存在信息获取滞后、决策依据不足、管理效率偏低等问题。智慧水利通过信息技术的深度应用，实现了管理方式的根本性变革。

\subsubsection{1.2.1
信息获取方式的变革}\label{ux4fe1ux606fux83b7ux53d6ux65b9ux5f0fux7684ux53d8ux9769}

\textbf{传统水利：} - 主要依靠人工观测和定期巡检 -
信息采集周期长，实时性差 - 监测点稀少，覆盖面有限 -
数据质量受人为因素影响较大

\textbf{智慧水利：} - 自动化、智能化感知设备全天候监测 -
实时数据采集和传输 - 监测点密集分布，覆盖面广 - 数据质量高，误差小

\subsubsection{1.2.2
数据处理能力的提升}\label{ux6570ux636eux5904ux7406ux80fdux529bux7684ux63d0ux5347}

\textbf{传统水利：} - 主要采用统计分析方法 - 处理数据量有限 -
分析深度不够，难以发现深层规律 - 预测精度有待提高

\textbf{智慧水利：} - 运用大数据技术处理海量信息 -
多维度、多层次数据分析 - 深度挖掘数据价值，发现隐藏规律 -
基于AI算法提升预测精度

\subsubsection{1.2.3
决策支持机制的优化}\label{ux51b3ux7b56ux652fux6301ux673aux5236ux7684ux4f18ux5316}

\textbf{传统水利：} - 主要依靠专家经验和历史类比 - 决策过程主观性较强 -
方案评估手段有限 - 应急响应速度慢

\textbf{智慧水利：} - 基于数据分析和模型计算的科学决策 -
客观量化的决策依据 - 多方案对比和效果预评估 - 快速响应和动态调整

\subsection{1.3
智慧水利发展历程}\label{ux667aux6167ux6c34ux5229ux53d1ux5c55ux5386ux7a0b}

\subsubsection{1.3.1
初期探索阶段（2009-2012年）}\label{ux521dux671fux63a2ux7d22ux9636ux6bb52009-2012ux5e74}

这一阶段主要是水利信息化的起步期，重点关注基础设施建设和数据采集能力提升。

\textbf{主要特征：} - 水文监测站网建设逐步完善 -
开始尝试自动化监测设备的应用 - 建立了一批水利数据库和信息系统 -
信息化应用主要集中在数据管理和简单查询

\textbf{代表性成果：} - 国家防汛抗旱指挥系统初步建成 -
水文信息采集传输系统网络化水平提升 - 水利工程安全监测系统开始规模化应用

\subsubsection{1.3.2
快速发展阶段（2013-2018年）}\label{ux5febux901fux53d1ux5c55ux9636ux6bb52013-2018ux5e74}

随着云计算、大数据技术的成熟，水利信息化进入快速发展期，应用范围和深度显著提升。

\textbf{主要特征：} - 云计算平台在水利行业开始应用 -
大数据分析技术逐步引入水利业务 - 移动互联网技术在水利管理中普及 -
跨部门、跨区域的信息共享机制初步建立

\textbf{代表性成果：} - 水利云平台建设启动 - 水利大数据中心建设推进 -
移动水利应用广泛部署 - 水利一张图建设取得重要进展

\subsubsection{1.3.3
智能化升级阶段（2019年至今）}\label{ux667aux80fdux5316ux5347ux7ea7ux9636ux6bb52019ux5e74ux81f3ux4eca}

进入新时代，人工智能、5G、数字孪生等前沿技术在水利行业深度应用，智慧水利建设全面提速。

\textbf{主要特征：} - 人工智能技术在水利预报预警中应用 -
数字孪生技术推动水利工程数字化转型 - 5G技术支撑水利万物互联 -
智慧水利标准体系逐步完善

\textbf{代表性成果：} - 数字孪生流域建设试点启动 -
AI驱动的洪水预报系统投入使用 - 5G+智慧水利应用示范项目实施 -
智慧水利总体框架和标准规范发布

\subsection{1.4
国家政策导向与行业发展现状}\label{ux56fdux5bb6ux653fux7b56ux5bfcux5411ux4e0eux884cux4e1aux53d1ux5c55ux73b0ux72b6}

\subsubsection{1.4.1
国家政策支持}\label{ux56fdux5bb6ux653fux7b56ux652fux6301}

近年来，国家高度重视智慧水利建设，出台了一系列政策文件，为智慧水利发展提供了有力的政策保障。

\textbf{重要政策文件：}

\begin{enumerate}
\def\labelenumi{\arabic{enumi}.}
\tightlist
\item
  \textbf{《``十四五''水安全保障规划》（2021年）}

  \begin{itemize}
  \tightlist
  \item
    明确提出构建智慧水利体系
  \item
    强调数字孪生流域建设
  \item
    推进水利新型基础设施建设
  \end{itemize}
\item
  \textbf{《关于推进智慧水利建设的指导意见》（2022年）}

  \begin{itemize}
  \tightlist
  \item
    确立了智慧水利建设的总体目标和重点任务
  \item
    提出了''需求牵引、应用至上、数字赋能、提升能力''的基本原则
  \item
    明确了数字孪生流域、数字孪生水网、数字孪生工程的建设内容
  \end{itemize}
\item
  \textbf{《数字中国建设整体布局规划》（2023年）}

  \begin{itemize}
  \tightlist
  \item
    将智慧水利作为数字中国建设的重要组成部分
  \item
    强调数字技术与水利业务的深度融合
  \item
    推动水利治理数字化转型
  \end{itemize}
\end{enumerate}

\subsubsection{1.4.2
行业发展现状}\label{ux884cux4e1aux53d1ux5c55ux73b0ux72b6}

\textbf{发展成果：} -
\textbf{基础设施日益完善}：建成了覆盖全国的水文监测站网，监测点数量达到12万余个
- \textbf{数据资源不断丰富}：累积了海量的水文、气象、工程等数据资源 -
\textbf{应用系统逐步完善}：建成了国家防汛抗旱指挥系统、水资源管理系统等一批重要应用
- \textbf{标准规范逐步健全}：制定了水利信息化相关标准规范100余项

\textbf{存在挑战：} -
\textbf{数据壁垒}：部门间、区域间数据共享仍存在障碍 -
\textbf{技术融合}：新技术与水利业务的深度融合有待加强 -
\textbf{人才缺乏}：既懂水利又精通信息技术的复合型人才不足 -
\textbf{标准统一}：行业标准的统一性和权威性需要提升

\subsection{1.5
智慧水利发展趋势}\label{ux667aux6167ux6c34ux5229ux53d1ux5c55ux8d8bux52bf}

\subsubsection{1.5.1
技术发展趋势}\label{ux6280ux672fux53d1ux5c55ux8d8bux52bf}

\begin{enumerate}
\def\labelenumi{\arabic{enumi}.}
\tightlist
\item
  \textbf{数字孪生技术将成为核心}

  \begin{itemize}
  \tightlist
  \item
    构建流域、水网、工程的数字孪生体
  \item
    实现物理世界与虚拟世界的实时交互
  \item
    支撑预报、预警、预演、预案功能
  \end{itemize}
\item
  \textbf{人工智能应用将更加深入}

  \begin{itemize}
  \tightlist
  \item
    AI算法在水文预报中的精度不断提升
  \item
    机器学习在工程安全监测中的应用扩展
  \item
    智能决策支持系统日趋成熟
  \end{itemize}
\item
  \textbf{云边端协同架构将成为主流}

  \begin{itemize}
  \tightlist
  \item
    云计算提供强大的计算和存储能力
  \item
    边缘计算实现数据的就近处理
  \item
    终端设备智能化水平不断提升
  \end{itemize}
\end{enumerate}

\subsubsection{1.5.2
应用发展趋势}\label{ux5e94ux7528ux53d1ux5c55ux8d8bux52bf}

\begin{enumerate}
\def\labelenumi{\arabic{enumi}.}
\tightlist
\item
  \textbf{全流域一体化管理}

  \begin{itemize}
  \tightlist
  \item
    打破行政边界，实现流域统一管理
  \item
    上下游、左右岸协调联动
  \item
    多部门、多层级协同决策
  \end{itemize}
\item
  \textbf{精准化调度控制}

  \begin{itemize}
  \tightlist
  \item
    基于实时数据的动态调度
  \item
    多目标优化的智能决策
  \item
    精细化的水资源配置
  \end{itemize}
\item
  \textbf{主动式风险防控}

  \begin{itemize}
  \tightlist
  \item
    从被动应对向主动防控转变
  \item
    风险识别和预警能力显著提升
  \item
    应急响应更加快速高效
  \end{itemize}
\end{enumerate}

\subsubsection{1.5.3
产业发展趋势}\label{ux4ea7ux4e1aux53d1ux5c55ux8d8bux52bf}

\begin{enumerate}
\def\labelenumi{\arabic{enumi}.}
\tightlist
\item
  \textbf{产业生态更加完善}

  \begin{itemize}
  \tightlist
  \item
    形成完整的智慧水利产业链
  \item
    产学研用深度融合
  \item
    标准化体系不断完善
  \end{itemize}
\item
  \textbf{市场规模快速增长}

  \begin{itemize}
  \tightlist
  \item
    智慧水利市场需求旺盛
  \item
    技术创新驱动产业发展
  \item
    国际合作与交流增强
  \end{itemize}
\item
  \textbf{商业模式不断创新}

  \begin{itemize}
  \tightlist
  \item
    平台化服务模式兴起
  \item
    数据价值化程度提升
  \item
    投融资模式多样化
  \end{itemize}
\end{enumerate}

\subsection{小结}\label{ux5c0fux7ed3}

智慧水利作为新时代水利发展的重要方向，通过新一代信息技术与传统水利的深度融合，实现了水利管理方式的根本性变革。从信息获取的自动化智能化，到数据处理的大数据化智能化，再到决策支持的科学化精准化，智慧水利为水利现代化发展注入了强大动力。

当前，我国智慧水利建设正处于重要战略机遇期，在国家政策的大力支持下，技术创新不断突破，应用场景日益丰富，产业生态逐步完善。面向未来，数字孪生、人工智能、云边端协同等技术将推动智慧水利向更高水平发展，为保障国家水安全、支撑经济社会高质量发展提供有力支撑。

作为水利专业的学生，深入理解智慧水利的发展脉络和技术特点，对于把握行业发展趋势、提升专业素养具有重要意义。在后续学习中，我们将进一步探讨智慧水利平台的技术架构和核心功能，为掌握智慧水利平台开发技能奠定坚实基础。

\input{chapters/chapter01/section01-02}  
\section{第三节
平台主要功能介绍}\label{ux7b2cux4e09ux8282-ux5e73ux53f0ux4e3bux8981ux529fux80fdux4ecbux7ecd}

\subsection{1.10
智慧水利平台功能体系}\label{ux667aux6167ux6c34ux5229ux5e73ux53f0ux529fux80fdux4f53ux7cfb}

智慧水利平台的功能体系是根据水利业务需求设计的，覆盖水利工作的各个方面。从功能定位上，可以将平台功能分为四大类：监测预警、调度决策、工程管理和信息服务。

\subsubsection{1.10.1
功能体系框架}\label{ux529fux80fdux4f53ux7cfbux6846ux67b6}

智慧水利平台的功能体系框架如下：

\begin{figure}
\centering
\includesvg{../../assets/images/chapter01_function_framework.svg}
\caption{智慧水利平台功能体系框架图}
\end{figure}

\textbf{监测预警功能}：实现对水利要素的全面感知和预警\\
\textbf{调度决策功能}：支持水利工程的科学调度和决策\\
\textbf{工程管理功能}：实现水利工程全生命周期管理\\
\textbf{信息服务功能}：提供面向各类用户的信息服务

这四类功能相互关联、相互支撑，共同构成完整的智慧水利平台功能体系。

\subsubsection{1.10.2
功能设计原则}\label{ux529fux80fdux8bbeux8ba1ux539fux5219}

智慧水利平台的功能设计应遵循以下原则：

\textbf{业务驱动原则}： - 功能设计以水利业务需求为出发点 -
充分考虑水利工作实际流程和习惯 - 注重解决实际业务问题和痛点

\textbf{用户体验原则}： - 界面设计简洁明了，符合用户使用习惯 -
操作流程简单高效，降低使用门槛 - 提供个性化配置，适应不同用户需求

\textbf{技术先进原则}： - 采用先进成熟的技术，确保系统稳定可靠 -
运用新技术提升系统智能化水平 - 保持技术架构开放性，支持未来扩展

\textbf{安全可靠原则}： - 功能设计充分考虑安全性要求 -
关键功能提供备份和容错机制 - 确保系统在各种条件下稳定运行

\subsection{1.11
监测预警功能}\label{ux76d1ux6d4bux9884ux8b66ux529fux80fd}

监测预警功能是智慧水利平台的基础功能，通过各类监测设备和技术手段，实现对水利要素的全面感知和预警。

\subsubsection{1.11.1
水文气象监测}\label{ux6c34ux6587ux6c14ux8c61ux76d1ux6d4b}

水文气象监测功能负责收集和处理水文气象数据，包括：

\textbf{降雨监测}： -
雨量站实时监测：接入各类雨量站数据，实时监测降雨情况 -
雨量分布分析：基于监测数据和雷达数据，分析降雨空间分布 -
累计雨量统计：计算不同时段的累计雨量，支持自定义统计区间

\textbf{水位流量监测}： -
水位实时监测：接入水位站数据，实时监测河湖水位变化 -
流量计算分析：基于水位数据和率定关系，计算河道流量 -
水位流量过程线：展示水位和流量的时间变化过程

\textbf{气象要素监测}： - 温度湿度监测：监测气温和湿度的变化 -
风向风速监测：监测风向和风速的变化 -
气压蒸发监测：监测气压和蒸发量的变化

\subsubsection{1.11.2
水利工程监测}\label{ux6c34ux5229ux5de5ux7a0bux76d1ux6d4b}

水利工程监测功能负责监测水利工程的运行状态，包括：

\textbf{大坝安全监测}： - 变形监测：监测大坝位移、沉降等变形情况 -
渗流监测：监测大坝渗流量和渗流压力 -
应力应变监测：监测大坝内部应力和应变状态

\textbf{水库运行监测}： - 水位监测：监测水库水位变化 -
库容计算：根据水位计算水库蓄水量 -
出入库流量监测：监测水库入库和出库流量

\textbf{闸门运行监测}： - 开度监测：监测闸门开启高度 -
启闭状态监测：监测闸门启闭状态 -
机电设备状态监测：监测启闭机等设备运行状态

\subsubsection{1.11.3 水资源监测}\label{ux6c34ux8d44ux6e90ux76d1ux6d4b}

水资源监测功能负责监测水资源状况，包括：

\textbf{水量监测}： - 地表水监测：监测河湖水量变化 -
地下水监测：监测地下水位和储量变化 -
用水计量监测：监测各类用水户的取水量

\textbf{水质监测}： - 常规指标监测：监测pH、溶解氧、电导率等常规指标 -
污染物指标监测：监测COD、氨氮、重金属等污染物指标 -
生物指标监测：监测藻类、细菌等生物指标

\textbf{水生态监测}： - 水温监测：监测水体温度变化 -
透明度监测：监测水体透明度变化 - 生物多样性监测：监测水生生物多样性状况

\subsubsection{1.11.4
预警预报功能}\label{ux9884ux8b66ux9884ux62a5ux529fux80fd}

预警预报功能负责基于监测数据进行分析预测和预警，包括：

\textbf{洪水预报}： - 降雨预报：基于气象数据预报未来降雨 -
洪水预报：预报河道洪峰流量和到达时间 -
水库调度预报：预报水库调度后的下游河道水情

\textbf{干旱预警}： - 干旱监测：监测土壤墒情和地表水缺乏程度 -
干旱预警：基于监测数据和气象预报发布干旱预警 -
供水保障分析：分析干旱条件下的供水保障能力

\textbf{水质预警}： - 水质异常监测：监测水质参数的异常变化 -
水质预警：发现水质异常时发出预警 - 污染溯源分析：分析水质污染的可能来源

\textbf{工程安全预警}： - 监测数据分析：分析工程监测数据的变化趋势 -
安全状态评价：评价工程的安全状态 - 异常预警：发现异常状态时发出预警

\subsection{1.12
调度决策功能}\label{ux8c03ux5ea6ux51b3ux7b56ux529fux80fd}

调度决策功能是智慧水利平台的核心功能，通过数据分析和模型计算，支持水利工程的科学调度和决策。

\subsubsection{1.12.1 防洪调度}\label{ux9632ux6d2aux8c03ux5ea6}

防洪调度功能支持洪水期间的科学调度决策，包括：

\textbf{洪水调度方案生成}： -
情景模拟：模拟不同降雨和调度情景下的洪水演进 -
方案优化：根据防洪目标优化调度方案 -
调度过程管理：管理调度方案的制定和执行过程

\textbf{梯级水库联合调度}： -
梯级水库信息整合：整合梯级水库的水情和工情信息 -
联合调度模型：建立梯级水库联合调度模型 -
优化调度计算：计算最优的联合调度方案

\textbf{风险分析与评估}： - 洪水风险评估：评估洪水风险等级和影响范围 -
调度风险评估：评估不同调度方案的风险 -
风险预警：对高风险区域和情况进行预警

\subsubsection{1.12.2 水资源调度}\label{ux6c34ux8d44ux6e90ux8c03ux5ea6}

水资源调度功能支持水资源的优化配置和利用，包括：

\textbf{供水调度}： - 需水分析：分析不同用户的需水量 -
可供水分析：分析可供水资源量 - 供水方案生成：生成最优的供水调度方案

\textbf{灌溉调度}： - 灌区需水计算：计算灌区需水量 -
灌溉制度优化：优化灌溉用水时间和数量 -
节水灌溉控制：实现精准灌溉和节水控制

\textbf{水电调度}： - 发电效益分析：分析水电站发电效益 -
发电调度优化：优化水电站的发电调度 -
电力系统协调：与电力系统协调优化调度

\subsubsection{1.12.3 应急调度}\label{ux5e94ux6025ux8c03ux5ea6}

应急调度功能支持突发事件下的应急处置，包括：

\textbf{应急预案管理}： - 预案制定：制定各类突发事件的应急预案 -
预案演练：开展应急预案的演练 - 预案评估：评估应急预案的有效性

\textbf{应急响应}： - 事件监测：监测突发事件的发生和发展 -
影响评估：评估突发事件的影响范围和程度 -
应急指挥：协调各方力量开展应急处置

\textbf{应急调度决策}： - 情景模拟：模拟突发事件下的不同处置情景 -
应急措施生成：生成应急处置措施 - 决策支持：为应急指挥提供决策支持

\subsubsection{1.12.4
决策支持系统}\label{ux51b3ux7b56ux652fux6301ux7cfbux7edf}

决策支持系统为水利调度决策提供全面支持，包括：

\textbf{数据分析}： - 多源数据融合：整合多源数据进行综合分析 -
趋势分析：分析数据变化趋势 - 关联分析：分析不同要素间的关联关系

\textbf{模型计算}： - 水文模型：降雨-径流模型、洪水演进模型等 -
水力模型：河道水力学模型、水库调度模型等 -
优化模型：多目标优化模型、决策树模型等

\textbf{方案评估}： - 效益评估：评估调度方案的经济效益 -
风险评估：评估调度方案的风险 - 综合评价：对调度方案进行综合评价

\subsection{1.13
工程管理功能}\label{ux5de5ux7a0bux7ba1ux7406ux529fux80fd}

工程管理功能支持水利工程全生命周期的科学管理，包括建设管理、运行管理和维护管理。

\subsubsection{1.13.1
工程建设管理}\label{ux5de5ux7a0bux5efaux8bbeux7ba1ux7406}

工程建设管理功能支持水利工程建设过程的管理，包括：

\textbf{项目管理}： - 计划管理：制定和管理工程建设计划 -
进度管理：监控和管理工程建设进度 - 质量管理：监控和管理工程建设质量

\textbf{设计管理}： - 设计文档管理：管理工程设计文档 -
设计审查：支持设计方案的审查 - 设计变更：管理设计变更过程

\textbf{施工管理}： - 施工监控：实时监控施工现场 -
工程量管理：管理工程量的计量和结算 - 验收管理：支持工程验收过程管理

\subsubsection{1.13.2
工程运行管理}\label{ux5de5ux7a0bux8fd0ux884cux7ba1ux7406}

工程运行管理功能支持水利工程的日常运行管理，包括：

\textbf{运行监控}： - 实时监控：实时监控工程运行状态 -
运行参数管理：管理工程运行参数 - 运行记录：记录工程运行情况

\textbf{调度管理}： - 调度计划管理：管理工程调度计划 -
调度指令下达：下达工程调度指令 - 调度执行监控：监控调度指令的执行情况

\textbf{运行评价}： - 运行效率评价：评价工程运行效率 -
运行成本分析：分析工程运行成本 - 运行绩效评价：评价工程运行绩效

\subsubsection{1.13.3
工程维护管理}\label{ux5de5ux7a0bux7ef4ux62a4ux7ba1ux7406}

工程维护管理功能支持水利工程的维护管理，包括：

\textbf{维修养护管理}： - 养护计划管理：制定和管理养护计划 -
维修工单管理：管理维修工单的创建和处理 - 养护记录管理：记录维修养护活动

\textbf{设备管理}： - 设备台账管理：管理设备的基本信息 -
设备状态监控：监控设备的运行状态 - 设备维护管理：管理设备的维护活动

\textbf{安全评价}： - 安全检查：支持工程安全检查 -
隐患管理：管理工程安全隐患 - 安全评估：评估工程的安全状况

\subsubsection{1.13.4 资产管理}\label{ux8d44ux4ea7ux7ba1ux7406}

资产管理功能支持水利工程资产的全生命周期管理，包括：

\textbf{资产台账管理}： - 基础信息管理：管理资产的基本信息 -
资产分类管理：按类型管理资产 - 资产变动管理：管理资产的增减变动

\textbf{资产评估}： - 资产价值评估：评估资产的价值 -
资产状况评估：评估资产的使用状况 - 资产效益评估：评估资产的使用效益

\textbf{资产决策支持}： - 资产规划：支持资产配置规划 -
投资决策：支持资产投资决策 - 处置决策：支持资产处置决策

\subsection{1.14
信息服务功能}\label{ux4fe1ux606fux670dux52a1ux529fux80fd}

信息服务功能负责向各类用户提供水利信息服务，包括管理决策服务、政务服务和公众服务。

\subsubsection{1.14.1
管理决策服务}\label{ux7ba1ux7406ux51b3ux7b56ux670dux52a1}

管理决策服务面向水利管理者和决策者，提供决策支持服务，包括：

\textbf{综合分析}： - 数据可视化：直观展示各类水利数据 -
统计分析：对水利数据进行统计分析 - 专题分析：针对特定主题进行深入分析

\textbf{形势研判}： - 水情分析：分析当前水情状况 -
发展趋势：预测水情发展趋势 - 风险预警：对潜在风险进行预警

\textbf{辅助决策}： - 情景模拟：模拟不同决策情景的效果 -
方案比较：比较不同决策方案的优劣 - 建议生成：生成决策建议

\subsubsection{1.14.2 政务服务}\label{ux653fux52a1ux670dux52a1}

政务服务面向政府部门和企事业单位，提供水利政务服务，包括：

\textbf{在线办事}： - 取水许可：提供取水许可申请和审批服务 -
工程审批：提供水利工程审批服务 - 水土保持：提供水土保持方案审批服务

\textbf{信息公开}： - 政策法规：公开水利政策法规信息 -
规划计划：公开水利规划和计划信息 - 统计数据：公开水利统计数据

\textbf{互动交流}： - 意见征集：征集公众对水利工作的意见 -
投诉举报：处理水利相关投诉和举报 - 咨询答复：回答水利相关咨询

\subsubsection{1.14.3 公众服务}\label{ux516cux4f17ux670dux52a1}

公众服务面向社会公众，提供水利信息和服务，包括：

\textbf{水情信息}： - 实时水情：提供实时水情信息 -
水质信息：提供水质信息 - 预警信息：提供洪水、干旱等预警信息

\textbf{科普宣传}： - 水利知识：普及水利科学知识 -
防灾避险：宣传防灾减灾知识 - 节水护水：宣传节水和水环境保护知识

\textbf{互动服务}： - 信息查询：提供水利信息查询服务 -
问题反馈：接收公众问题反馈 - 智能问答：提供智能问答服务

\subsubsection{1.14.4
移动应用服务}\label{ux79fbux52a8ux5e94ux7528ux670dux52a1}

移动应用服务通过移动终端提供随时随地的水利信息服务，包括：

\textbf{移动监测}： - 现场数据采集：支持现场巡查和数据采集 -
移动监控：随时查看监测数据和视频 - 异常上报：及时上报发现的异常情况

\textbf{移动办公}： - 移动审批：随时处理审批事项 -
移动会议：支持远程视频会议 - 消息通知：实时接收重要消息通知

\textbf{移动服务}： - 位置服务：提供基于位置的水利信息 -
实时查询：随时查询水利信息 - 互动交流：便捷参与水利工作互动

\subsection{习题与思考}\label{ux4e60ux9898ux4e0eux601dux8003}

\begin{enumerate}
\def\labelenumi{\arabic{enumi}.}
\tightlist
\item
  分析智慧水利平台的功能体系框架，说明各类功能之间的关系。
\item
  监测预警是智慧水利平台的基础功能，请详细描述水文气象监测的主要内容和技术手段。
\item
  调度决策是智慧水利平台的核心功能，请分析防洪调度功能的主要内容和实现方式。
\item
  工程管理功能如何支持水利工程的全生命周期管理？请结合实例进行说明。
\item
  针对不同类型的用户，智慧水利平台应该提供哪些信息服务？请设计一套合理的信息服务方案。
\end{enumerate}

\subsection{参考文献}\label{ux53c2ux8003ux6587ux732e}

\begin{enumerate}
\def\labelenumi{\arabic{enumi}.}
\tightlist
\item
  水利部. 水利信息化标准体系{[}Z{]}. 2020.
\item
  王浩, 王建华, 等. 智慧水利理论体系与技术路线{[}J{]}. 水利学报, 2021,
  52(3): 1-10.
\item
  陈守煜, 张国新, 等. 智慧水利信息服务关键技术研究{[}J{]}. 水利信息化,
  2019(4): 10-16.
\item
  刘昌明, 王中根, 等. 水文水资源管理信息系统{[}M{]}. 北京：科学出版社,
  2017.
\item
  尚松浩, 李致家, 等. 大数据环境下水利决策支持系统构建{[}J{]}. 中国水利,
  2020(11): 15-19.
\end{enumerate}
